% =====================================================================
%  Section V - Results (HARBox).  Tables: sections/results_tables.tex, rows generated by
%  scripts/tools/paper_results_tables.py from results/harbox and results/harbox_fedchar.
%  Numbers quoted in the text are taken from those generated tables (user split 0 unless stated).
% =====================================================================
\section{Results}\label{sec:results}

% =====================================================================
%  Result tables (rows generated by scripts/tools/paper_results_tables.py).
%  Included by sections/results.tex.
% =====================================================================

\begin{table}[!t]
\centering
\caption{HARBox under the published protocol (Protocol A: random 72/8/20 split of 50--500 windows
per user, all 120 users, 100 rounds, half of the users per round). Mean test accuracy over users and
of the worst-served $10\,\%$. Published numbers are copied from FedCHAR~\cite{li2023fedchar} (benign
setting); \emph{ours} are re-runs under the same protocol. Sorted by accuracy.}
\label{tab:sota}
\footnotesize
\setlength{\tabcolsep}{3pt}
\renewcommand{\arraystretch}{1.08}
\begin{tabularx}{\columnwidth}{@{}Y l r r l@{}}
\toprule
Method & Model & Acc. & Worst-10\,\% & Source \\
\midrule
\inputrows{generated/res_tab_sota.tex}
\bottomrule
\end{tabularx}
\end{table}

\begin{table*}[!t]
\centering
\caption{HARBox, leakage-free Protocol B (84 clients, per-activity chronological split with a
one-window gap). Each cell: mean user accuracy / pooled macro-F1 on the clients' test parts.
Methods with a shared classifier are scored after every client fine-tuned that classifier on its own
training part (``+FT''), so every row compares personalized models. Best per encoder in bold. Budgets differ between encoders
(\cref{tab:hparams}), so compare methods within a column.}
\label{tab:main}
\footnotesize
\setlength{\tabcolsep}{5pt}
\renewcommand{\arraystretch}{1.1}
\begin{tabular}{@{}l c c c c c@{}}
\toprule
\inputrows{generated/res_tab_main.tex}
\bottomrule
\end{tabular}
\end{table*}

\begin{table}[!t]
\centering
\caption{New users (24 held-out users, Protocol B): mean per-user macro-F1 with zero, $1$, $5$ and
$10$ labelled windows per activity. ``transd.'': the same V5 model with the transductive
zero-shot rule (\cref{eq:refine}), which also initializes the few-shot head.}
\label{tab:newusers}
\footnotesize
\setlength{\tabcolsep}{3.5pt}
\renewcommand{\arraystretch}{1.08}
\begin{tabular}{@{}l l r r r r@{}}
\toprule
Encoder & Model & 0 & 1 & 5 & 10 \\
\midrule
\inputrows{generated/res_tab_newusers.tex}
\bottomrule
\end{tabular}
\end{table}

\begin{table}[!t]
\centering
\caption{Component ablation of \method{} (LIMU-BERT, Protocol B, one component removed at a
time). Acc.: mean user accuracy; $\Delta$: macro-F1 change against the complete V5; new users:
mean per-user macro-F1 with $0$ and $10$ labelled windows per activity.}
\label{tab:ablation}
\footnotesize
\setlength{\tabcolsep}{2.4pt}
\renewcommand{\arraystretch}{1.08}
\begin{tabularx}{\columnwidth}{@{}Y r r r r r r@{}}
\toprule
Variant & Acc. & F1 & $\Delta$ & W-10\,\% & 0 & 10 \\
\midrule
\inputrows{generated/res_tab_ablation.tex}
\bottomrule
\end{tabularx}
\end{table}

\begin{table}[!t]
\centering
\caption{Client cost of full fine-tuning (V0) and \method{} (V5) on HARBox: uplink per client update,
peak training memory (MB) and time per local update (s).}
\label{tab:cost}
\footnotesize
\setlength{\tabcolsep}{2.6pt}
\renewcommand{\arraystretch}{1.08}
\begin{tabular}{@{}l r r r r r r r@{}}
\toprule
 & \multicolumn{3}{c}{Uplink} & \multicolumn{2}{c}{Memory} & \multicolumn{2}{c}{Time} \\
Encoder & V0 & V5 & Ratio & V0 & V5 & V0 & V5 \\
\midrule
\inputrows{generated/res_tab_cost.tex}
\bottomrule
\end{tabular}
\end{table}

\begin{table*}[!t]
\centering
\caption{Repeated user splits (Protocol B, seeds $0$--$2$; $n$: number of splits): mean\,$\pm$\,standard
deviation. Enrolled clients: mean user accuracy, macro-F1 and the mean accuracy of the worst-served
$10\,\%$ (shared classifiers after ``+FT''). New users: mean per-user macro-F1 zero-shot (V5:
transductive rule) and with $10$ labelled windows per activity.}
\label{tab:seeds}
\footnotesize
\setlength{\tabcolsep}{4pt}
\renewcommand{\arraystretch}{1.1}
\begin{tabular}{@{}l l c c c c c c@{}}
\toprule
 & & & \multicolumn{3}{c}{Enrolled clients} & \multicolumn{2}{c}{New users} \\
\cmidrule(lr){4-6}\cmidrule(lr){7-8}
Method & Model & $n$ & Acc. & Macro-F1 & Worst-10\,\% & 0-shot & 10-shot \\
\midrule
\inputrows{generated/res_tab_seeds.tex}
\bottomrule
\end{tabular}
\end{table*}


% ---------------------------------------------------------------------
\subsection{Comparison with Published Results}\label{sec:results:sota}

\Cref{tab:sota} places our models next to the published HARBox results under the published
protocol. Our re-run of the published baseline, FedAvg with the 300-unit MLP, reaches a mean user
accuracy of $\ResAFedAvgMLPPlain$ ($\ResAFedAvgMLP$ after local fine-tuning of its classifier),
against $0.694$ reported by FedCHAR. The difference comes from the training recipe shared by all
our methods (AdamW, augmentation, label-skew handling), so the comparison below is conservative for
the baselines. Every personalized model we trained exceeds the best published result (FedCHAR,
$0.917$), without its hierarchical clustering. Ditto, the method on which FedCHAR builds, reaches
$\ResADittoMLP$ on the MLP. \method{} on LIMU-BERT reaches $\ResAVfive$ ($\ResAVfiveAll$ with all-layer
LoRA). Ditto started from the same pretrained LIMU-BERT reaches $\ResADittoFT$. The same table exposes
the weakness of the protocol: \emph{purely local} training, with no federation at all, already reaches
$0.878$, and centralized training on the pooled data ($0.881$--$0.896$) stays below the personalized
models. With $\HBLeakA\,\%$ of the test windows overlapping a training window, Protocol~A largely
measures how well a model memorizes a user's recording. A personal copy of every weight (Ditto) is the
best memorizer, and its lead over \method{} ($0.009$) disappears under the leakage-free Protocol~B
(\cref{sec:results:seeds}). We therefore report Protocol~A for comparability only and base all
conclusions on Protocol~B.

% ---------------------------------------------------------------------
\subsection{Leakage-Free Benchmark}\label{sec:results:main}

\Cref{tab:main} reports the main grid under Protocol~B.

\emph{\method{} is the best of the methods built on a pretrained encoder (V0--V5) on every encoder.}
It reaches a mean user accuracy of $0.861$ on LIMU-BERT, $0.875$ on ssl-wearables and $0.812$ on HALO,
which is $0.068$, $0.065$ and $0.098$ above full federated fine-tuning (V0). This holds although V0 is
scored after every client fine-tuned its classifier locally (``+FT''). On UniMTS, the slowest encoder,
with only $40$ rounds of $4$ local steps, the margin shrinks to $0.001$ accuracy and $0.004$ macro-F1
over V3, but the worst-served $10\,\%$ of the clients still gain ($0.370$ vs.\ $0.346$ accuracy). On LIMU-BERT
and ssl-wearables, V5 also exceeds the centralized models trained on the pooled data of all clients
($0.833$ and $0.847$ at best), and on HALO it matches the pretrained centralized model ($0.813$). Pooling
the data removes the communication constraint but not the heterogeneity between users. A model that is
personal to every user addresses the heterogeneity.

\emph{Federated LoRA without personalization.} With a shared classifier, federated LoRA (V2, V3)
matches or beats full fine-tuning on the two convolutional encoders (macro-F1 $0.833$ vs.\ $0.823$ on
ssl-wearables, $0.823$ vs.\ $0.817$ on UniMTS) but not on the two Transformers ($0.786$ vs.\ $0.813$ on
LIMU-BERT, $0.727$ vs.\ $0.738$ on HALO). Naive factor averaging (V2) and exact aggregation (V3) are
within $0.006$ macro-F1 on every encoder: at these ranks and this number of local steps the cross
terms of \cref{eq:naive-bias} are small. Heterogeneous ranks (V4) cost $0.006$--$0.024$ macro-F1
against homogeneous rank $8$ (V3).

\emph{LoRA placement matters more than the aggregation rule.} LIMU-BERT shares one Transformer block
across its four layers, so adapting $W_q$ and $W_v$ gives only $2.3$\,k LoRA parameters. Wrapping all
six linear layers of the block ($8.1$\,k parameters) raises V3 from $0.786$ to $0.817$ macro-F1, above
full fine-tuning ($0.813$), and V5 from $0.863$ to $0.873$. The uplink grows from $26$ to $39$\,kB per
update, still $5.2\times$ less than full fine-tuning. Raising the rank instead helps less: on $W_q,W_v$,
V3 gains steadily from rank $2$ to $16$ ($0.769$, $0.782$, $0.786$, $0.798$ macro-F1), but rank $16$
($4.6$\,k parameters) stays $0.019$ below all-layer rank $8$.

\emph{The strongest baseline.} Ditto started from the pretrained LIMU-BERT keeps a personal copy of the
whole model on every client. On split~$0$ it reaches $0.871$ accuracy and $0.877$ macro-F1, and
\method{} with all-layer LoRA matches it ($0.870$, $0.873$). Over three splits (\cref{tab:seeds}),
\method{} is level or ahead on every enrolled-client metric ($0.877$ vs.\ $0.875$ accuracy, $0.883$ vs.\
$0.882$ macro-F1, $0.581$ vs.\ $0.569$ worst-$10\,\%$), with $5.2\times$ less uplink and one local update per
round instead of two. On the MLP, Ditto ($0.868$) is likewise far ahead of FedAvg ($0.766$) and of local
training ($0.824$): under a leakage-free split, personalization helps, and so does federation.

% ---------------------------------------------------------------------
\subsection{Robustness Across User Splits}\label{sec:results:seeds}

\Cref{tab:seeds} repeats the LIMU-BERT and MLP experiments on two further user splits (a new
assignment of users to clients, held-out and pool users each time). The ranking is stable. \method{}
reaches $0.868\pm0.012$ mean user accuracy, against $0.784\pm0.013$ for V0 and $0.769\pm0.012$ for
FedAvg: its gain of $0.084$ over V0 is about seven standard deviations of the split-to-split
variation. With all-layer LoRA, \method{} reaches $0.877\pm0.015$, level with Ditto on the pretrained
encoder ($0.875\pm0.008$); with $10$ labelled windows per activity, it is the best model for new users
($0.852$ vs.\ $0.827$--$0.845$). The new-user scores vary much more between splits ($\pm0.04$--$0.06$
zero-shot), because each split has only $24$ held-out users.

% ---------------------------------------------------------------------
\subsection{Fairness}\label{sec:results:fairness}

Personalization matters most for the worst-served users. Averaged over the three splits, the mean
accuracy of the $10\,\%$ worst clients is $0.275$--$0.328$ for the shared LIMU-BERT models (FedAvg, V0,
V3) and $0.284$ for FedAvg on the MLP, against $0.568$ for \method{}, $0.567$ for Ditto and $0.551$ for
local training. The worst clients of the shared models score far below chance on their test parts
(e.g., $0.09$ for one user). This is consistent with self-labelling errors or unusual phone placement in
crowd-sourced data, which a shared classifier cannot absorb but a personal model can.

% ---------------------------------------------------------------------
\subsection{New Users}\label{sec:results:newusers}

\Cref{tab:newusers} reports the held-out users, who took part in no training.

\emph{Zero-shot}, the shared classifier of full fine-tuning transfers better than \method{}'s
nearest-prototype rule on every encoder: $0.688$ vs.\ $0.591$ on LIMU-BERT, $0.797$ vs.\ $0.775$ on
ssl-wearables, $0.760$ vs.\ $0.715$ on UniMTS and $0.568$ vs.\ $0.534$ on HALO.
The embedding of \method{} is trained jointly with personal classifiers and local norms, which a new
user does not have yet. Three measures narrow the gap without labels. Initializing a new user with the
clients' average local parameters (\cref{sec:method:inference}) raises zero-shot from $0.575$ to
$0.591$. The transductive rule (\cref{eq:refine}) raises it further to $0.626$ ($0.699\pm0.063$ over the
three splits, against $0.743\pm0.048$ for V0), UniMTS's from $0.715$ to $0.732$ and HALO's from
$0.534$ to $0.554$. On ssl-wearables, where AdaBN has already adapted every normalization layer to the
user, it adds almost nothing ($0.775$ to $0.776$). All-layer LoRA, together with the transductive rule, reaches $0.665$ on LIMU-BERT
($0.727\pm0.054$ over three splits), within $0.02$ of V0 ($0.743$) and of Ditto on the pretrained
encoder ($0.745$).

\emph{With labels}, the picture changes. With $10$ labelled windows per activity ($20$\,s),
\method{} reaches $0.854$ on LIMU-BERT, against $0.829$ for V0 and $0.807$ for V3. On HALO it is the best
model at every labelled budget ($0.764$, $0.817$ and $0.815$ with $1$, $5$ and $10$ windows). On the
two BatchNorm encoders, the shared models remain ahead with labels ($0.872$--$0.874$ vs.\ $0.860$ on
ssl-wearables, $0.859$ vs.\ $0.827$ on UniMTS, whose few-shot budget is cut to $30$ steps). Keeping the
AdaBN statistics fixed during few-shot adaptation (\cref{sec:method:inference}) matters here: when a
batch of a few labelled windows was allowed to overwrite them, the $1$-shot score of ssl-wearables fell
below its zero-shot score ($0.751$); with the fix, it is $0.796$.

% ---------------------------------------------------------------------
\subsection{Communication, Memory and Time}\label{sec:results:cost}

\Cref{tab:cost} compares the client cost of \method{} with full fine-tuning. The uplink per update
shrinks $8\times$ on LIMU-BERT ($204$ to $26$\,kB), $52\times$ on ssl-wearables ($42.1$\,MB to
$811$\,kB), $38\times$ on UniMTS ($20.9$\,MB to $555$\,kB) and $95\times$ on HALO ($18.5$\,MB to
$194$\,kB). Peak training memory falls by $34\,\%$ on UniMTS and $24\,\%$ on HALO, where the optimizer
state of the full model matters, and by $10\,\%$ on LIMU-BERT. On ssl-wearables, the V5 run peaked
higher ($480$ vs.\ $389$\,MB). Memory on the convolutional encoders includes cuDNN workspaces, which
vary between runs, and activations dominate for every encoder, because LoRA still back-propagates
through the whole network. The time per local update does not separate the methods reliably: the
runs shared one GPU with other experiments, and the same configuration varied by up to $\pm50\,\%$
between runs. Over a round, compute time is small next to the transfer of a full model on a mobile link.

% ---------------------------------------------------------------------
\subsection{Ablation}\label{sec:results:ablation}

\Cref{tab:ablation} removes or adds one component of \method{} at a time on LIMU-BERT, with LoRA on
$W_q,W_v$ as in the main grid.

\emph{Local normalization and input adapter} are the most important component. Without them,
macro-F1 drops by $0.025$ and the worst-$10\,\%$ accuracy from $0.533$ to $0.461$. Without all local
parts and prototypes (V4), macro-F1 drops by $0.093$.

\emph{The personal classifier} adds little once the norms are local. A shared classifier that every
client fine-tunes before testing reaches the same accuracy ($0.862$ vs.\ $0.861$). \emph{Prototype
alignment} does not change the enrolled clients ($-0.001$) or zero-shot ($+0.005$), but it raises the
$10$-shot score by $0.010$. \emph{Device tiers} cost nothing: a single rank $8$ for every device gives
the same macro-F1 ($+0.003$), so the low-rank devices of \method{} are not penalized.

\emph{Augmentation} is a trade-off. Without it, the enrolled clients gain $0.011$ macro-F1, but new
users lose $0.010$ zero-shot and $0.009$ with $10$ labels. We keep it, because the paper targets new
users as well as enrolled ones. \emph{Averaged local parts for new users} raise every new-user score
($+0.016$ zero-shot, $+0.011$ at $10$ shots) for a one-time upload of $2$\,kB per client.

\emph{Negative results.} Three further ideas did not help. An average of the clients' personal
classifiers as the zero-shot rule changed nothing ($+0.002$). A second, federated classifier for new
users (dual head) cost $0.009$ macro-F1 for a $+0.007$ zero-shot gain. A private copy of the LoRA
adapters, fine-tuned in the last personal update (the PEFT counterpart of Ditto's personal model), did
not improve the enrolled clients either with $20$ or with $60$ steps ($\le+0.003$). By contrast, all-layer
LoRA ($+0.011$ macro-F1, $+0.044$ zero-shot) and the transductive zero-shot rule ($+0.035$) are clear gains.

% ---------------------------------------------------------------------
\subsection{Discussion and Limitations}\label{sec:results:limits}

\emph{Evaluation protocol.} The largest effect in this study is not a method but the protocol:
under the random split of the published results, purely local training approaches the best published
federated method, while under a leakage-free split every model loses several points and federation
clearly beats local training. We recommend that future HARBox results report the per-activity
chronological split (or another split that keeps overlapping windows on one side) alongside any
random split.

\emph{New users versus enrolled users.} Local normalization and personal classifiers are what make
\method{} strong for enrolled users, and they are also why its zero-shot accuracy for brand-new users
is lower than that of a shared model. Averaged local parts and the transductive rule narrow this gap, and
a few labelled windows close it on LIMU-BERT and HALO. Deployments that must serve many
unlabelled newcomers may combine \method{} with a shared fallback model.

\emph{Personalization depth.} Ditto personalizes every weight and pays with a full-model uplink and two
local updates per round. \method{} personalizes only the norms, the input adapter and the classifier,
on top of shared low-rank adapters, and reaches the same accuracy once the adapters cover the whole
Transformer block. A private copy of the adapters adds nothing, which suggests that on HARBox the
user-specific part of the signal lives in the input statistics and the decision rule rather than in the
encoder weights.

\emph{Downlink.} Exact aggregation changes the frozen weights, and our accounting conservatively
charges their full dense size at every download. For the convolutional encoders, this makes the
downlink comparable to full fine-tuning. A factored encoding of the residual (\cref{sec:method:cost}) is
left for future work.

\emph{Scope.} The large encoders train with reduced budgets on our single 6\,GB GPU, so absolute numbers
across encoders are not budget-matched. HALO is our re-implementation, pretrained only on the
unlabelled windows of 12 pool users. LIMU-BERT and the MLP baselines are evaluated on three user splits;
the other encoders on one.
